% ============================================================
% Sección 11: Patrones de diseño agénticos
% ============================================================
\section{Patrones de diseño agénticos}

\begin{frame}{Cuatro patrones recurrentes}
  \begin{enumerate}
    \item \textbf{Reflection} (auto-crítica)
    \item \textbf{Planning} (planeación)
    \item \textbf{Tool use} (uso de herramientas)
    \item \textbf{Multi-agent collaboration} (colaboración multiagente)
  \end{enumerate}
  \vfill
  \footnotesize Popularizados por Andrew Ng / DeepLearning.AI como los
  bloques básicos con los que se construye casi cualquier sistema agéntico.
\end{frame}

\begin{frame}{Patrón: Reflection}
  \centering
  \begin{tikzpicture}[node distance=1.4cm]
    \node[cajaLLM] (gen) {Generar respuesta/borrador};
    \node[cajaAccento, right=2.2cm of gen] (crit) {Autocriticar\\(el mismo u otro LLM)};
    \node[cajaAgente, below=of gen] (fin) {Versión mejorada};
    \draw[flecha] (gen) -- (crit);
    \draw[flecha] (crit) -- node[right, etiquetaFlecha]{feedback} (fin);
    \draw[flechaPunteada] (fin.west) to[out=180,in=180, looseness=1.2] (gen.west);
  \end{tikzpicture}
  \vfill
  \footnotesize El agente revisa su propio trabajo antes de entregarlo —
  como un borrador y su edición.
\end{frame}

\begin{frame}{Patrón: Planning}
  \centering
  \resizebox{!}{0.58\textheight}{%
  \begin{tikzpicture}[node distance=0.7cm]
    \node[cajaAgente] (obj) {Objetivo complejo};
    \node[cajaLLM, below=of obj] (plan) {Descomponer en subtareas};
    \node[cajaHerramienta, below=of plan] (s1) {Subtarea 1 $\to$ 2 $\to$ 3 \dots};
    \node[cajaAccento, below=of s1] (fin) {Ejecutar y combinar};
    \draw[flecha] (obj) -- (plan);
    \draw[flecha] (plan) -- (s1);
    \draw[flecha] (s1) -- (fin);
  \end{tikzpicture}}
  \vfill
  \footnotesize El agente traza un plan \emph{antes} de actuar, en vez de
  improvisar acción por acción.
\end{frame}

\begin{frame}{Patrón: Tool use}
  \centering
  \begin{tikzpicture}[node distance=1.6cm]
    \node[cajaLLM] (llm) {LLM};
    \node[cajaHerramienta, above right=0.8cm and 2cm of llm] (t1) {Calculadora};
    \node[cajaHerramienta, right=2cm of llm] (t2) {Buscador};
    \node[cajaHerramienta, below right=0.8cm and 2cm of llm] (t3) {Ejecutor de código};
    \draw[flechaDoble] (llm) -- (t1);
    \draw[flechaDoble] (llm) -- (t2);
    \draw[flechaDoble] (llm) -- (t3);
  \end{tikzpicture}
  \vfill
  \footnotesize El patrón que ya vimos en RAG y ReAct: el LLM decide
  \emph{cuándo} y \emph{con qué argumentos} invocar cada herramienta.
\end{frame}

\begin{frame}{Patrón: Multi-agent collaboration}
  \centering
  \footnotesize Ya visto en la sección anterior: orquestador-trabajadores,
  pipeline, debate.
  \vfill
  \begin{itemize}
    \item Se combina con los otros tres: cada agente puede tener su propio
      ciclo de \emph{planning + tool use + reflection}.
    \item Es el patrón de \textbf{mayor nivel}: compone a los demás.
  \end{itemize}
\end{frame}

\begin{frame}{Los cuatro patrones combinados}
  \centering
  \resizebox{0.8\textwidth}{0.48\textheight}{%
  \begin{tikzpicture}[node distance=0.9cm]
    \node[cajaBase, fill=colGris!15, minimum width=8.5cm] (top) {\textbf{Multi-agent collaboration}};
    \node[cajaAgente, below=of top, minimum width=2.5cm] (p1) {Planning};
    \node[cajaHerramienta, right=0.4cm of p1, minimum width=2.5cm] (p2) {Tool use};
    \node[cajaAccento, right=0.4cm of p2, minimum width=2.5cm] (p3) {Reflection};
    \draw[flecha] (top.south) -- ++(0,-0.3) -| (p1.north);
    \draw[flecha] (top.south) -- (p2.north);
    \draw[flecha] (top.south) -- ++(0,-0.3) -| (p3.north);
  \end{tikzpicture}}
  \vfill
  \footnotesize En la práctica, un buen agente combina varios de estos
  patrones a la vez, no uno solo.
\end{frame}
